% Report for the CISC3024 project. All numbers come from numbers.tex (generated from results/*.json by scripts/build_report_assets.py).
% Red [STUDENT: ...] marks are things only the student can supply or must confirm; remove them (and set \showmarks to 0) before upload.
\documentclass[10pt]{article}
\usepackage[a4paper,margin=1.75cm,top=1.6cm,bottom=1.6cm]{geometry}
\usepackage{amsmath,amssymb,booktabs,graphicx,xcolor,url,enumitem,microtype,caption}
\usepackage[colorlinks=true,linkcolor=blue!50!black,urlcolor=blue!50!black,citecolor=blue!50!black]{hyperref}
\setlist{nosep,leftmargin=1.2em}
\captionsetup{font=small,labelfont=bf,skip=3pt}
\setlength{\parskip}{2pt}
\renewcommand{\arraystretch}{1.05}
% GENERATED by scripts/build_report_assets.py from results/*.json - do not edit
\newcommand{\officialRows}{0 & 0.972 & 46.57 & 7.3 & 0.796 & 0.419 & 0.00\% & 15 \\
101 & 0.975 & 46.84 & 7.3 & 0.808 & 0.416 & 0.00\% & 15 \\
202 & 0.970 & 46.47 & 7.1 & 0.801 & 0.433 & 0.00\% & 15 \\
303 & 0.973 & 47.05 & 6.8 & 0.748 & 0.445 & 0.00\% & 15 \\
404 & 0.964 & 46.79 & 7.0 & 0.780 & 0.425 & 0.00\% & 15 \\}
\newcommand{\chkDet}{yes}
\newcommand{\chkSeed}{yes}
\newcommand{\chkMixed}{0.96}
\newcommand{\chkSingle}{0.95}
\newcommand{\manifestRows}{8{,}029}
\newcommand{\params}{5{,}098{,}915}
\newcommand{\sizeMB}{19.8}
\newcommand{\sampleSpeed}{18}
\newcommand{\trainHours}{1.3}
\newcommand{\nSeeds}{61}
\newcommand{\nFifteen}{60}
\newcommand{\nBelow}{1}
\newcommand{\recallMean}{0.433}
\newcommand{\recallMin}{0.395}
\newcommand{\recallMax}{0.488}
\newcommand{\fidMean}{46.72}
\newcommand{\fidMin}{46.14}
\newcommand{\fidMax}{47.35}
\newcommand{\accMean}{0.971}
\newcommand{\accMin}{0.962}
\newcommand{\accMax}{0.979}
\newcommand{\precMean}{0.785}
\newcommand{\precMin}{0.746}
\newcommand{\precMax}{0.814}
\newcommand{\recallSD}{0.019}
\newcommand{\memMax}{0.00}
\newcommand{\lowSeed}{31415}
\newcommand{\lowRecall}{0.3949}
\newcommand{\lowZ}{2.0}
\newcommand{\pFifteen}{98.4}
\newcommand{\pFifteenLo}{91.3}
\newcommand{\pFifteenHi}{99.7}
\newcommand{\wideAtFifteen}{12}
\newcommand{\wideRecall}{0.416}
\newcommand{\widePrec}{0.829}
\newcommand{\wideAcc}{0.979}
\newcommand{\smallAtFifteen}{15}
\newcommand{\smallRecall}{0.433}
\newcommand{\smallPrec}{0.788}
\newcommand{\smallAcc}{0.972}
\newcommand{\progressRows}{v009 & 1{,}399 & 12 / -- / -- & 0.936 & 60.6 & 0.333 & 0.759 & 4.35\% \\
v010 & 1{,}929 & 13 / -- / -- & 0.942 & 58.0 & 0.319 & 0.769 & 1.65\% \\
v010 + mirror & 1{,}929 & 14 / 13 / 13 & 0.954 & 50.6 & 0.388 & 0.768 & 0.00\% \\
v011 & 5{,}536 & 15 / 14 / 15 & 0.981 & 48.5 & 0.401 & 0.819 & 0.00\% \\
v012 & 4{,}610 & 14 / 14 / 15 & 0.976 & 48.0 & 0.381 & 0.828 & 0.00\% \\
v013 & 4{,}922 & 14 / -- / -- & 0.975 & 49.0 & 0.353 & 0.837 & 0.00\% \\
v014 (final) & 8{,}029 & 15 / 15 / 15 & 0.972 & 46.6 & 0.423 & 0.802 & 0.00\% \\}
\newcommand{\auditGood}{160}
\newcommand{\auditPct}{80.0}
\newcommand{\auditCI}{73.9--85.0}
\newcommand{\auditCat}{63/75}
\newcommand{\auditCatPct}{84.0}
\newcommand{\auditDog}{97/125}
\newcommand{\auditDogPct}{77.6}
\newcommand{\auditPersonal}{43/49}
\newcommand{\auditAI}{117/151}
\newcommand{\auditEars}{33}
\newcommand{\auditBad}{40}
\newcommand{\auditRecrop}{7/18}

\newcommand{\showmarks}{1}
\newcommand{\student}[1]{\ifnum\showmarks=1\textcolor{red}{[STUDENT: #1]}\fi}
\newcommand{\zipurl}{https://github.com/Nvb-flipped/REPO-NAME/releases/download/final-submission/project_submission.zip}
\title{\vspace{-0.9cm}\textbf{Cat and dog head generation from a self-curated pool}\vspace{-0.1cm}}
\author{Sou Chak Tong \quad Student ID: U-C3-2652-3 \quad CISC3024 Pattern Recognition Project}
\date{}
\begin{document}
\maketitle\vspace{-0.9cm}

\section*{Summary}
I curated a training set from the 76{,}250-image LAION pool and trained a class-conditional diffusion model from scratch on it.
The final training set has \textbf{8{,}029 crops} (3{,}096 cat, 4{,}933 dog): 1{,}929 I reviewed myself and 6{,}100 found by AI-assisted screening and approved in an agent visual review.
The generator is a \params-parameter pixel-space U-Net (\sizeMB\,MB), trained for 80k updates with mirroring, EMA and classifier-free guidance, sampled with 50 deterministic DDIM steps.
On the unchanged official evaluator, seed 0 gives accuracy 0.972, FID 46.57, KID $7.3{\times}10^{-3}$, precision 0.796, recall 0.419 and 0.00\,\% memorization, i.e.\ \textbf{15/15 marks}.
Over \nSeeds{} generation seeds, \nFifteen{} score 15/15 ($\pFifteen\,\%$, Wilson 95\,\% interval $\pFifteenLo$--$\pFifteenHi\,\%$); seed \lowSeed{} (one I chose myself) scores 14/15 because its recall is \lowRecall{}, just below 0.40.
Recall is therefore the one metric with little margin. My personal audit of 200 training crops found \auditGood{} good (\auditPct\,\%, 95\,\% interval \auditCI\,\%); the main error is ears cut by the square.
Visual quality is weaker than the metrics suggest: by my own count 39 of 128 samples in the seed-0 grid are unrecognizable (Section~3).
\textbf{Submission ZIP:} \expandafter\url\expandafter{\zipurl}\ \student{replace by the real link; add the OneDrive backup link if used}.

\section{Data}
\textbf{Pipeline.} Grounding DINO Tiny~\cite{liu2023groundingdino} proposed head boxes for all 76{,}250 pool images; I turned them into square crops with the official geometry (60--80\,\% head occupancy, no padding, no rotation) and added bounded re-squaring where the box fell outside the image or covered only the face. A CLIP~\cite{radford2021clip} screen and a targeted second detector pass recovered images with no detection, giving \textbf{63{,}115 candidate crops}. Qwen3.5-9B~\cite{qwen35} (served with vLLM~\cite{kwon2023vllm}) screened every candidate (6{,}718 passed). A CLIP linear probe trained on my own verdicts ranked the rest so that review went best-first. About 17{,}000 crops then received an \emph{agent visual review} against the exact handout rules (real photo; exactly one centred, roughly upright head with ears; 60--80\,\% of the side; $\ge$32 px wide); 7{,}083 were judged valid. Same-photo overlaps, near-duplicates of my own crops (pixel distance, dHash, source reposts) and holdout sources were removed (983 crops). Together with my 1{,}929 personally reviewed crops (50 batches) this gives the final \textbf{8{,}029}. The evaluator's classifier, Inception network and all dev images and features were never used for curation. The agent's verdicts are not my own; the audit below is.

\begin{table}[h]
\centering\small\setlength{\tabcolsep}{5pt}
\begin{tabular}{lr}
\toprule
Step & Images \\
\midrule
Raw pool (LAION photographs, mixed with other animals, objects and scenes) & 76{,}250 \\
Candidate head crops (detector squares + bounded geometry rescue + targeted rescue of detector misses) & 63{,}115 \\
Passed the Qwen3.5-9B screen (a looser second pass only ordered the review) & 6{,}718 \\
Agent visual review under the handout rules: crops reviewed / judged valid & 17{,}174 / 7{,}083 \\
After removing same-photo overlaps, near-duplicates of my crops, duplicate families and holdout sources & 6{,}100 \\
Added: crops I reviewed personally (50 batches) & 1{,}929 \\
\midrule
\textbf{Final training manifest} (3{,}096 cat, 4{,}933 dog; 7{,}719 distinct pool images) & \textbf{\manifestRows} \\
Pool-only holdout, never trained on (31 cat, 32 dog) & 63 \\
\bottomrule
\end{tabular}
\caption{Data pipeline with the number of images after each step (the review also covered crops the screen rejected, so the steps are not a strict funnel). The official evaluator rebuilt all \manifestRows{} crops from the manifest exactly.}
\end{table}

\textbf{Audit.} I labelled the 200 crops drawn by the official \texttt{audit.py} from my student ID without seeing any agent verdict: \textbf{\auditGood{} of 200 good = \auditPct\,\%} (Wilson 95\,\% interval \auditCI\,\%). Cats \auditCat{} (\auditCatPct\,\%), dogs \auditDog{} (\auditDogPct\,\%); my own reviewed rows \auditPersonal{}, AI-found rows \auditAI{}. Of the \auditBad{} bad crops, \auditEars{} have ears cut by the square, 4 a head under 60\,\%, 2 a covered head or ears and 1 a cut-off head; none was a cartoon, a second head or a wrong label. Ear clipping was worst in re-cropped squares (\auditRecrop{} audited re-crops). The cause is that detector boxes often cover only the face and the agent judged ear margins on small thumbnails. A residual rate of about 3\,\% non-photos (painting, figurine, print on a mug) was found in a 96-crop agent sample and removed. A missed-head sample of 200 unflagged sources found about 1\,\% real heads (Wilson 0.3--3.6\,\%). \textbf{Next:} re-review every crop at 300 px or larger with an explicit ear-margin check, widen or drop failures, retrain and redraw an audit; I did not do this because it needs a full re-review, retraining and a new audit and would void the current one.

\begin{figure}[h]
\centering\includegraphics[width=0.94\linewidth]{fig_audit.pdf}
\caption{Result of my personal audit of 200 crops: share judged good by group (dot, with 95\,\% Wilson interval) and the reasons for the 40 bad crops. The intervals overlap, so my own rows and the AI-found rows cannot be reliably told apart.}
\end{figure}

\section{Methods}
\textbf{Model.} A U-Net~\cite{ronneberger2015unet} on $3{\times}32{\times}32$ images scaled to $[-1,1]$ with four resolutions (32, 16, 8, 4), widths 32/64/96/128, two residual blocks per level, GroupNorm~\cite{wu2018groupnorm} (8 groups), SiLU and dropout 0.1. Self-attention (4 heads) acts at 8$\times$8, 4$\times$4, the bottleneck and, as an addition that helped in my ablations, at 16$\times$16 in the encoder. A conditioning vector, an MLP of the sinusoidal time embedding plus a learned class embedding with three entries (cat, dog, null), gives a scale and shift to every residual block. \params{} parameters, randomly initialized, no pretrained weights.

\textbf{Diffusion and training.} With the cosine schedule~\cite{nichol2021improved}, $\bar\alpha_t=f(t)/f(0)$, $f(t)=\cos^2\!\big(\tfrac{t/T+s}{1+s}\tfrac{\pi}{2}\big)$, $T{=}1000$, $s{=}0.008$, a crop $x_0$ is noised as $x_t=\sqrt{\bar\alpha_t}\,x_0+\sqrt{1-\bar\alpha_t}\,\epsilon$, $\epsilon\sim\mathcal N(0,I)$ \cite{ho2020ddpm}. The network predicts the noise, trained with
\begin{equation}
\mathcal L(\theta)=\mathbb E_{x_0,y,t,\epsilon}\,\big\|\epsilon_\theta(x_t,t,\tilde y)-\epsilon\big\|_2^2,\qquad \tilde y=\varnothing\ \text{with probability }0.1,\ \text{else } y,
\end{equation}
$t$ uniform on $\{0,\dots,T{-}1\}$ and each exact manifest crop mirrored left--right with probability 0.5 (the evaluator already counts mirrors as copies). AdamW~\cite{loshchilov2019adamw} with learning rate $2{\times}10^{-4}$ after 2{,}000 warm-up updates, weight decay $10^{-4}$, batch 128, gradient clipping 1.0, bfloat16, and an exponential moving average $\theta_{\text{EMA}}\leftarrow0.9995\,\theta_{\text{EMA}}+0.0005\,\theta$ whose weights are submitted.

\textbf{Sampling.} DDIM~\cite{song2020ddim} with $S{=}50$ uniformly spaced steps and classifier-free guidance~\cite{ho2022cfg} $\tilde\epsilon=\epsilon_\varnothing+w(\epsilon_y-\epsilon_\varnothing)$, $w{=}2.75$, plus perturbed-attention guidance~\cite{ahn2024pag} $\tilde\epsilon\mathrel{+}=0.15(\epsilon_y-\epsilon_y^{\text{pert}})$ for $149\le t\le949$ (one encoder attention layer replaced by identity), a temporal rescale $\tilde\epsilon\leftarrow s_t\tilde\epsilon$ with $s_t=\frac{\bar\alpha_t\sigma^2+1-\bar\alpha_t}{\bar\alpha_t\sigma^2/1.03+1-\bar\alpha_t}$, $\sigma{=}0.5$, class-embedding noise up to 10\,\% at high noise, and FreeU~\cite{si2024freeu} strength 2.0 in the decoder. Each step forms $\hat x_0=\mathrm{clip}_{[-1,1]}\big((x_t-\sqrt{1-\bar\alpha_t}\,\tilde\epsilon)/\sqrt{\bar\alpha_t}\big)$, $\hat\epsilon=(x_t-\sqrt{\bar\alpha_t}\hat x_0)/\sqrt{1-\bar\alpha_t}$ and $x_{t'}=\sqrt{\bar\alpha_{t'}}\hat x_0+\sqrt{1-\bar\alpha_{t'}}\hat\epsilon$. The starting noise of each output comes from a private CPU generator seeded by the caller's seed (one draw per output position), the caller's random state is preserved, and a fixed physical batch of 32 keeps every output independent of the other labels in a batch.

\section{Results}
\begin{table}[t]
\centering\small\setlength{\tabcolsep}{4.5pt}
\begin{tabular}{rrrrrrrr}
\toprule
Seed & Accuracy & FID & KID ($10^{-3}$) & Precision & Recall & Memorization & Marks \\
\midrule
\officialRows
\midrule
\multicolumn{8}{l}{All \nSeeds{} seeds: accuracy \accMean{} (\accMin--\accMax), FID \fidMean{} (\fidMin--\fidMax), recall \recallMean{} $\pm$ \recallSD{} (\recallMin--\recallMax),}\\
\multicolumn{8}{l}{precision \precMean{} (\precMin--\precMax), memorization at most \memMax\,\%; \nFifteen{} of \nSeeds{} seeds earn 15 marks.}\\
\bottomrule
\end{tabular}
\caption{Official evaluator (1{,}000 images per class, batch 500, grading mode) run on \texttt{submission/} itself; seed 0 is the headline. Thresholds for 5 marks each: accuracy $\ge$0.90, FID $\le$50, recall $\ge$0.40; memorization must stay $\le$5\,\%. \params{} parameters, \sizeMB\,MB on disk. Hardware: one NVIDIA RTX 5090 (shared) and an Intel Core Ultra 7 265K. Training: \trainHours\,h for the 80k-update checkpoint (two runs shared the GPU; the full 160k-update run took 2.6\,h). Sampling: about \sampleSpeed\,s per 1{,}000 images on the GPU.}
\end{table}

\begin{table}[t]
\centering\small\setlength{\tabcolsep}{4pt}
\begin{tabular}{lrcrrrrr}
\toprule
Model & Images & Marks (seeds 0/101/202) & Acc. & FID & Recall & Prec. & Mem. \\
\midrule
\progressRows
\bottomrule
\end{tabular}
\caption{Development path with the same evaluator; metrics are means over the listed seeds. v009 is the 12/15 model I started this phase with; v010 adds all my reviewed batches, then mirroring; v011--v014 add AI-found images (v012 and v013 use stricter agent review, so fewer images). Guidance and sampler settings differ slightly between rows (v009/v010: guidance 3.25, v010 + mirror: 3.0, v011--v014: 2.75).}
\end{table}

\begin{figure}[t]
\centering\includegraphics[width=0.80\linewidth]{fig_grid_seed0.png}
\caption{The evaluator's own sample grid of the submitted model, seed 0: rows 1--4 are cats, rows 5--8 are dogs (16 images per row). The image file is saved by the unchanged official evaluator.}
\end{figure}

\begin{figure}[t]
\centering\includegraphics[width=0.92\linewidth]{fig_progress.pdf}
\caption{More valid training images (right) lift recall and lower FID. The strictly reviewed v012 and v013 sets are smaller and have lower recall than v011.}
\end{figure}

\begin{figure}[t]
\centering
\begin{minipage}[c]{0.50\linewidth}\centering
\includegraphics[width=\linewidth]{fig_seeds.pdf}
\captionof{figure}{Recall of all \nSeeds{} evaluated seeds, sorted. Only seed \lowSeed{} falls below 0.40.}
\end{minipage}\hfill
\begin{minipage}[c]{0.46\linewidth}\centering\small
\begin{tabular}{@{}lr@{}}
\toprule
Hard limit or check & Result \\
\midrule
Parameters ($\le$20\,M) & \params \\
Size on disk ($\le$100\,MB) & \sizeMB\,MB \\
Memorization ($\le$5\,\%) & 0.00\,\% \\
Deterministic generation & \chkDet \\
Different seed, different images & \chkSeed \\
Accuracy, single-class vs mixed batches & \chkSingle{} vs \chkMixed \\
Manifest rebuilt by the evaluator & \manifestRows{} of \manifestRows \\
\bottomrule
\end{tabular}
\captionof{table}{Hard limits and interface checks reported by the evaluator (seed 0).}
\end{minipage}
\end{figure}

\textbf{What drove the result.} Mirroring was the largest single change (v010: 13 to 14 marks, memorization 1.65\,\% to 0\,\%), then more valid training images: recall rose from about 0.32 with 1.9k images to 0.42 with 8k, and FID fell below 50 (Fig.~2). On the final data a wider 10.7M-parameter model raised accuracy (\wideAcc{} vs \smallAcc) and precision (\widePrec{} vs \smallPrec) but lowered mean recall (\wideRecall{} vs \smallRecall) and earned 15 marks on only \wideAtFifteen{} of 15 seeds (the final model: \smallAtFifteen{} of 15), so I kept the small model.

\textbf{Sample quality and limits.} I counted unrecognizable samples in the official grids, i.e.\ tiles where I can hardly or not at all tell whether they show a cat, a dog or even an animal because of distortion or misplaced head parts: \textbf{39 of 128} in seed 0 (cats 10/64, dogs 29/64) and \textbf{60 of 128} in seed 31415 (cats 21/64, dogs 39/64). Dogs fail more often than cats, and the official metrics do not show this: accuracy is 0.97 and precision 0.80. The agent had described the grid as having ``about 6--8 distorted dog tiles'', which my count shows was far too optimistic. Other limits: recall margin (seed \lowSeed{}), ear clipping in the data, only 8k training images, and the hidden test set and seed may behave differently from the dev set. The memorization check passed on every seed (0.00\,\%), but the evaluator compares only against training crops and mirrors with a pixel metric.

\section{What did not work}
Full details for all 94 experiment folders are in \texttt{EXPERIMENTS.md}.
\begin{itemize}
\item \textbf{Memorization at about 1.2k images} was the binding limit early on: 13 of 94 experiment folders have 0 marks because memorization exceeded 5\,\% (up to 74\,\%).
\item \textbf{Alternative objectives and architectures} at matched budget (Min-SNR~\cite{hang2023minsnr}, EDM~\cite{karras2022edm}, velocity prediction, pixel patches, flow matching, Haar wavelets) were all worse than the plain noise-prediction U-Net (FID 87--220 versus 72); self-conditioning~\cite{chen2023analogbits} failed to train usefully; dropout 0.15/0.25 and higher weight decay did not help.
\item \textbf{Sampler tricks} (FreeU, PAG, temporal rescale, class-noise annealing) each gained a few FID points and are in the final sampler, but alone could not get below FID 60 at 1.4k images. DPM-Solver++~\cite{lu2022dpmpp}, stochastic DDIM, CFG rescale and interval guidance~\cite{kynkaanniemi2024interval} gave no reliable gain.
\item \textbf{Cleaner but smaller data lowered recall}: my first strict agent review (v012, v013) kept fewer valid heads, and recall fell from 0.40 to 0.35--0.38. I then switched to the exact handout rules.
\item \textbf{Longer training and a wider network} did not help the final model: the 120k checkpoint lost a mark on 1 of 15 seeds, and the 10.7M-parameter model scored 15 on only 9--12 of 15 seeds because its recall was lower.
\end{itemize}

\section{Working with coding agents}
\paragraph{What I handed over and what I decided.} I used two coding agents as my main development tools: Codex on my first PC (30 October to 19 November) and Claude Code on a Linux GPU PC (from 19 November). I gave them the project requirements, my goals and an execution plan, and delegated most of the implementation: the data curation pipeline, the DDPM model, the training experiments and the collection of evaluation results, and later the search of the whole pool, containerized GPU jobs, the tests and the release tooling. My role was mainly to decide the direction and to judge whether the results were reliable. I chose a conditional diffusion model with classifier-free guidance, and decided to expand the training set when FID and recall were still below the target. I also wanted to recover as many valid cat and dog heads from the pool as possible (I estimated 9{,}000--10{,}000) instead of only the minimum. For the data I reviewed the proposed crops myself, corrected labels and crop boundaries, and decided which images to keep during duplicate review. In the last days I let the agent choose model and sampler settings within a fixed time window and reviewed the outcome afterwards; the final checkpoint was then chosen by a rule fixed before the results were known.

\paragraph{Where the agents helped.} They were most useful for repetitive, time-consuming work. They scanned the 76{,}250 pool images, detected candidate heads (63{,}115 candidates in the end), prepared crops and built a review interface that made selection much easier. At first I inspected every image and typed measurements and checklist items by hand. After several batches I asked the agent to pre-fill its suggested labels, measurements and rejection reasons, to calculate occupancy automatically and to change how the checklist behaves when I adjust a crop, so that I could concentrate on checking images and fixing mistakes. They also automated training, checkpoint evaluation and comparisons, which let me study many alternatives (model sizes, training length, augmentation, sampling settings) without managing each run: the project has 94 experiment folders and 210 tests.

\paragraph{How I checked their work.} My agent log records three checks. \emph{Data:} the agent reported that its approved crops had only about 3\,\% problems; I labelled 200 random crops against the handout without seeing its verdicts and found 20\,\% bad, 33 of 40 with cut ears. \emph{Results:} the agent reported 15/15 on every one of 60 seeds; I reran the unchanged official evaluator on \texttt{submission/}, and my own seed \lowSeed{} scored 14/15 (recall \lowRecall). \emph{Interface:} I wrote six probes (same seed twice, seed 5 against 6, image 0 with the other labels changed, one label against batches of 4, 40 and 65, shape and range, preserved random state), all exactly identical, and a hand calculation of one noising and recovery example matched \texttt{diffusion.py} (difference $5.6{\times}10^{-17}$); I found no mistake there. Earlier, on the first PC, I also reviewed 200 proposed crops against their source images (accepting 14 and rejecting 186, while correcting species labels and crop boundaries), reviewed 200 sources the detector had rejected and found 23 eligible heads it had missed, and compared 30 source-and-crop pairs, re-measuring borderline cases, for example a crop I enlarged from 250 to 260 pixels to bring the estimated head occupancy to about 76.8\,\%. These checks taught me that code which runs correctly can still produce proposals and measurements that are wrong, so I could not trust a suggestion without checking it against the requirements.

\paragraph{Where they got things wrong or overstated, and what I changed.} The first detector often proposed wrong species labels, unsuitable crop boundaries or images without a valid head, and it missed some eligible heads; I corrected these in my reviews and asked for a workflow that made doubtful cases easier to verify. The audit above showed that the later AI-assisted review missed ear clipping, and the agent's claims of ``15/15 on every seed'' and of ``about 6--8 distorted dogs'' per grid were too optimistic (my count is 29 of 64 dogs in seed 0). I added AGENTS.md rules to check ear margins at larger size, to write ``X of Y seeds'' with the spread, and to count unrecognizable tiles exactly for cats and dogs separately. A different problem appeared with Qwen3.5-4B for AI-assisted cleaning. Its first results looked promising on images I had already reviewed, but a fresh sample from the whole candidate pool showed that roughly a quarter to a third of the accepted images were invalid: other species, several heads, unsuitable crops or non-photographs. The agent also found that its prompt contained a program-generated head-size hint that the model seemed to copy instead of judging the image. We removed the hint, added a larger vision-language model and further checks, and AI-curated images were only allowed into the training set after a blind sample of the pool confirmed their quality. An early agent review rule that was stricter than the handout cost recall (v012, v013), so it was replaced by the exact handout rules. On 14 November the agent also reported that it had scheduled work to resume after the usage limit reset, but the work had not started when expected; when I questioned this, it acknowledged that setting up a schedule does not mean the task ran.

Overall, the agents were very helpful for implementation and experimentation, especially with a large pool and many training runs. But a finished task or a convincing result does not mean it is correct: I still had to examine the actual outputs, question unexpected results and decide on evidence instead of accepting what the agent reported.

\label{mainend}
\clearpage
\bibliographystyle{plain}
\bibliography{references}
\appendix
\section*{Appendix A: settings}
\small
\begin{tabular}{@{}lp{0.80\linewidth}@{}}
\toprule
Item & Value \\
\midrule
Manifest & \texttt{submission/manifest.csv} (v014), 8{,}029 rows, holdout of 63 pool crops (31 cat, 32 dog) \\
Model & channels 32/64/96/128, 2 blocks per level, 4 heads, conditioning 256, dropout 0.1, extra 16$\times$16 encoder attention \\
Training & seed 17, batch 128 (2 micro-batches of 64), AdamW $2{\times}10^{-4}$, $\beta=(0.9,0.999)$, wd $10^{-4}$, warm-up 2{,}000, EMA 0.9995, mirror 0.5, null label 0.1, bf16 \\
Sampler & DDIM 50 steps, $w=2.75$, PAG 0.15, temporal rescale 1.03 ($\sigma=0.5$), class noise 0.1, FreeU 2.0, execution batch 32 \\
Selected checkpoint & 80{,}000 updates of a 160{,}000-update run; chosen among 80k--160k by a rule fixed before the results; see \texttt{DECISION\_LOG.md} \\
Reproduction & \texttt{README.md} lists one command per number; the numbers above are generated by \texttt{scripts/build\_report\_assets.py}. All experiments, commands and tests ran on Ubuntu Linux (one RTX 5090); generation and \texttt{demo.py} are written for CPU or NVIDIA GPU on any system but were tested on Linux only \\
\bottomrule
\end{tabular}
\end{document}
